% Options for packages loaded elsewhere
\PassOptionsToPackage{unicode}{hyperref}
\PassOptionsToPackage{hyphens}{url}
\PassOptionsToPackage{dvipsnames,svgnames,x11names}{xcolor}
%
\documentclass[
  11pt,
]{article}
\usepackage{amsmath,amssymb}
\usepackage{iftex}
\ifPDFTeX
  \usepackage[T1]{fontenc}
  \usepackage[utf8]{inputenc}
  \usepackage{textcomp} % provide euro and other symbols
\else % if luatex or xetex
  \usepackage{unicode-math} % this also loads fontspec
  \defaultfontfeatures{Scale=MatchLowercase}
  \defaultfontfeatures[\rmfamily]{Ligatures=TeX,Scale=1}
\fi
\usepackage{lmodern}
\ifPDFTeX\else
  % xetex/luatex font selection
\fi
% Use upquote if available, for straight quotes in verbatim environments
\IfFileExists{upquote.sty}{\usepackage{upquote}}{}
\IfFileExists{microtype.sty}{% use microtype if available
  \usepackage[]{microtype}
  \UseMicrotypeSet[protrusion]{basicmath} % disable protrusion for tt fonts
}{}
\makeatletter
\@ifundefined{KOMAClassName}{% if non-KOMA class
  \IfFileExists{parskip.sty}{%
    \usepackage{parskip}
  }{% else
    \setlength{\parindent}{0pt}
    \setlength{\parskip}{6pt plus 2pt minus 1pt}}
}{% if KOMA class
  \KOMAoptions{parskip=half}}
\makeatother
\usepackage{xcolor}
\usepackage[margin=2.2cm]{geometry}
\usepackage{longtable,booktabs,array}
\usepackage{calc} % for calculating minipage widths
% Correct order of tables after \paragraph or \subparagraph
\usepackage{etoolbox}
\makeatletter
\patchcmd\longtable{\par}{\if@noskipsec\mbox{}\fi\par}{}{}
\makeatother
% Allow footnotes in longtable head/foot
\IfFileExists{footnotehyper.sty}{\usepackage{footnotehyper}}{\usepackage{footnote}}
\makesavenoteenv{longtable}
\setlength{\emergencystretch}{3em} % prevent overfull lines
\providecommand{\tightlist}{%
  \setlength{\itemsep}{0pt}\setlength{\parskip}{0pt}}
\setcounter{secnumdepth}{5}
\usepackage{lmodern}
\usepackage[T1]{fontenc}
\usepackage{amsmath,amssymb}
\usepackage{newunicodechar}
\newunicodechar{}{\ensuremath{\bar\nu}}
\newunicodechar{}{\ensuremath{\hat\nu}}
\newunicodechar{}{\ensuremath{\tilde\sigma}}
\newunicodechar{}{\ensuremath{\bar v}}
\newunicodechar{}{\ensuremath{\hat\sigma}}
\newunicodechar{}{\ensuremath{\bar X}}
\newunicodechar{}{\ensuremath{\hat X}}
\newunicodechar{}{\ensuremath{\tilde W}}
\newunicodechar{}{\ensuremath{\dot\alpha}}
\newunicodechar{}{\ensuremath{\dot\beta}}
\newunicodechar{}{\ensuremath{\hat f}}
\newunicodechar{}{\ensuremath{\bar t}}
\newunicodechar{}{\ensuremath{\hat V}}
\newunicodechar{Ĉ}{\ensuremath{\hat C}}
\newunicodechar{ĉ}{\ensuremath{\hat c}}
\newunicodechar{ġ}{\ensuremath{\dot g}}
\newunicodechar{ᵀ}{\ensuremath{{}^{\mathsf{T}}}}
\newunicodechar{ℝ}{\ensuremath{\mathbb{R}}}
\newunicodechar{−}{\ensuremath{-}}
\newunicodechar{²}{\ensuremath{{}^{2}}}
\newunicodechar{³}{\ensuremath{{}^{3}}}
\newunicodechar{¹}{\ensuremath{{}^{1}}}
\newunicodechar{⁴}{\ensuremath{{}^{4}}}
\newunicodechar{⁸}{\ensuremath{{}^{8}}}
\newunicodechar{⁺}{\ensuremath{{}^{+}}}
\newunicodechar{⁻}{\ensuremath{{}^{-}}}
\newunicodechar{₂}{\ensuremath{{}_{2}}}
\newunicodechar{½}{\ensuremath{\tfrac12}}
\newunicodechar{¼}{\ensuremath{\tfrac14}}
\newunicodechar{·}{\ensuremath{\cdot}}
\newunicodechar{×}{\ensuremath{\times}}
\newunicodechar{÷}{\ensuremath{\div}}
\newunicodechar{±}{\ensuremath{\pm}}
\newunicodechar{σ}{\ensuremath{\sigma}}
\newunicodechar{φ}{\ensuremath{\varphi}}
\newunicodechar{ν}{\ensuremath{\nu}}
\newunicodechar{τ}{\ensuremath{\tau}}
\newunicodechar{π}{\ensuremath{\pi}}
\newunicodechar{ρ}{\ensuremath{\rho}}
\newunicodechar{λ}{\ensuremath{\lambda}}
\newunicodechar{μ}{\ensuremath{\mu}}
\newunicodechar{κ}{\ensuremath{\kappa}}
\newunicodechar{ω}{\ensuremath{\omega}}
\newunicodechar{β}{\ensuremath{\beta}}
\newunicodechar{α}{\ensuremath{\alpha}}
\newunicodechar{γ}{\ensuremath{\gamma}}
\newunicodechar{θ}{\ensuremath{\theta}}
\newunicodechar{η}{\ensuremath{\eta}}
\newunicodechar{χ}{\ensuremath{\chi}}
\newunicodechar{ψ}{\ensuremath{\psi}}
\newunicodechar{ξ}{\ensuremath{\xi}}
\newunicodechar{δ}{\ensuremath{\delta}}
\newunicodechar{ε}{\ensuremath{\varepsilon}}
\newunicodechar{Σ}{\ensuremath{\Sigma}}
\newunicodechar{Δ}{\ensuremath{\Delta}}
\newunicodechar{Λ}{\ensuremath{\Lambda}}
\newunicodechar{Π}{\ensuremath{\Pi}}
\newunicodechar{Ξ}{\ensuremath{\Xi}}
\newunicodechar{Γ}{\ensuremath{\Gamma}}
\newunicodechar{Φ}{\ensuremath{\Phi}}
\newunicodechar{Θ}{\ensuremath{\Theta}}
\newunicodechar{Ψ}{\ensuremath{\Psi}}
\newunicodechar{∫}{\ensuremath{\int}}
\newunicodechar{∂}{\ensuremath{\partial}}
\newunicodechar{√}{\ensuremath{\surd}}
\newunicodechar{∞}{\ensuremath{\infty}}
\newunicodechar{≈}{\ensuremath{\approx}}
\newunicodechar{≥}{\ensuremath{\ge}}
\newunicodechar{≤}{\ensuremath{\le}}
\newunicodechar{≠}{\ensuremath{\ne}}
\newunicodechar{≡}{\ensuremath{\equiv}}
\newunicodechar{→}{\ensuremath{\to}}
\newunicodechar{←}{\ensuremath{\leftarrow}}
\newunicodechar{⇒}{\ensuremath{\Rightarrow}}
\newunicodechar{⇔}{\ensuremath{\Leftrightarrow}}
\newunicodechar{↦}{\ensuremath{\mapsto}}
\newunicodechar{∈}{\ensuremath{\in}}
\newunicodechar{∝}{\ensuremath{\propto}}
\newunicodechar{∏}{\ensuremath{\prod}}
\newunicodechar{⊥}{\ensuremath{\perp}}
\newunicodechar{‖}{\ensuremath{\|}}
\newunicodechar{⟨}{\ensuremath{\langle}}
\newunicodechar{⟩}{\ensuremath{\rangle}}
\newunicodechar{∩}{\ensuremath{\cap}}
\newunicodechar{⊆}{\ensuremath{\subseteq}}
\newunicodechar{…}{\ensuremath{\ldots}}
\newunicodechar{̄}{}
\newunicodechar{̂}{}
\newunicodechar{̃}{}
\newunicodechar{̇}{}
\ifLuaTeX
  \usepackage{selnolig}  % disable illegal ligatures
\fi
\IfFileExists{bookmark.sty}{\usepackage{bookmark}}{\usepackage{hyperref}}
\IfFileExists{xurl.sty}{\usepackage{xurl}}{} % add URL line breaks if available
\urlstyle{same}
\hypersetup{
  pdftitle={COS Deep-Dive --- Priority Question Set},
  pdfauthor={Computational Finance (WI4151), TU Delft --- Prof.~Fang Fang},
  colorlinks=true,
  linkcolor={blue},
  filecolor={Maroon},
  citecolor={Blue},
  urlcolor={Blue},
  pdfcreator={LaTeX via pandoc}}

\title{COS Deep-Dive --- Priority Question Set}
\author{Computational Finance (WI4151), TU Delft --- Prof.~Fang Fang}
\date{}

\begin{document}
\maketitle

{
\hypersetup{linkcolor=blue}
\setcounter{tocdepth}{3}
\tableofcontents
}
Fang invented COS. If she asks one hard question, it is probably here.
Goal: I can do everything on the checklist \textbf{closed-book, on
paper, out loud}.

\hypertarget{mastery-checklist-tick-when-you-can-do-it-cold}{%
\section{Mastery checklist (tick when you can do it
cold)}\label{mastery-checklist-tick-when-you-can-do-it-cold}}

\begin{itemize}
\tightlist
\item[$\square$]
  Derive the COS pricing formula for a European option from the
  Fourier-cosine expansion of the density.
\item[$\square$]
  State where the characteristic function enters and why COS only needs
  the cf, not the density.
\item[$\square$]
  Explain why a \textbf{cosine} expansion (even extension, real
  coefficients) rather than a full complex Fourier series.
\item[$\square$]
  Write the payoff coefficients V\_k for a call/put and explain the
  closed form.
\item[$\square$]
  Choose the truncation range {[}a, b{]} from cumulants: a, b = c1 ±
  L·sqrt(c2 + sqrt(c4)). Explain each term.
\item[$\square$]
  Explain the convergence rate: exponential for smooth densities,
  algebraic otherwise --- and what kills it.
\item[$\square$]
  \textbf{My Assignment 2 bug:} why a small c2 (short maturity) makes
  {[}a, b{]} too narrow to cover the spot grid, and exactly how I fixed
  it.
\item[$\square$]
  COS for Heston: plug in the Heston cf; state the Feller condition.
\item[$\square$]
  COS for Barrier options (Lecture 10) --- how the method adapts.
\item[$\square$]
  COS for American options (Lecture 9) --- the approach used.
\item[$\square$]
  COS vs Carr--Madan FFT: when and why COS wins.
\end{itemize}

\hypertarget{high-probability-questions-have-claude-code-grill-me-on-these}{%
\section{High-probability questions (have Claude Code grill me on
these)}\label{high-probability-questions-have-claude-code-grill-me-on-these}}

\begin{enumerate}
\def\labelenumi{\arabic{enumi}.}
\item
  Derive the COS European call price. Where does the cf enter?

  \textbf{A.} Start from the discounted-expectation price in the
  log-state \(y\) (e.g.~\texttt{y\ =\ ln(S\_T/K)}):

\begin{verbatim}
P(x0,T) = e^{-rT} ∫_R V(y,T) f(y|x0) dy.
\end{verbatim}

  I do not know \(f\), but I know its characteristic function
  \texttt{φ\_\{X\_T\}(t;x0)\ =\ E{[}e\^{}\{i\ t\ X\_T\}\textbar{}x0{]}\ =\ ∫\ e\^{}\{ity\}\ f(y)\ dy},
  i.e.~\(φ\) is the Fourier transform of \(f\).

  \textbf{Step 1 --- truncate.} Pick \([a,b]\) (cumulant rule, Q2)
  containing essentially all the mass and restrict the integral to
  \([a,b]\).

  \textbf{Step 2 --- cosine-expand the density.} On \([a,b]\) write the
  half-range cosine series (prime halves \(k=0\)):

\begin{verbatim}
f(y) ≈ Σ'_{k≥0} A_k cos(kπ (y-a)/(b-a)),   A_k = (2/(b-a)) ∫_a^b f(y) cos(kπ(y-a)/(b-a)) dy.
\end{verbatim}

  \textbf{Step 3 --- A\_k from φ (the engine).} Write \(cos = Re{exp}\),
  pull out the phase \texttt{e\^{}\{-ikπa/(b-a)\}}, and extend the
  integral to \(R\) (the \emph{only} approximation --- range
  truncation):

\begin{verbatim}
A_k ≈ F_k = (2/(b-a)) Re{ φ(kπ/(b-a); x0) · exp(-ikπ a/(b-a)) }.
\end{verbatim}

  The remaining
  \texttt{∫\_R\ f(y)\ e\^{}\{i(kπ/(b-a))y\}\ dy\ =\ φ(kπ/(b-a))}
  exactly. No integral survives.

  \textbf{Step 4 --- payoff coefficients.} Define
  \texttt{V\_k\ =\ (2/(b-a))\ ∫\_a\^{}b\ V(y,T)\ cos(kπ(y-a)/(b-a))\ dy}.
  For the call \texttt{V\ =\ K(e\^{}y-1)\^{}+} lives on \(y≥0\), so

\begin{verbatim}
V_k^call = (2/(b-a)) K [ χ_k(0,b) − ψ_k(0,b) ].
\end{verbatim}

  \textbf{Step 5 --- assemble.} Insert the expansion, swap \(Σ\)/\(∫\)
  (uniform convergence of the cosine series for smooth \(f\)), and use
  \(V_k\):

\begin{verbatim}
P ≈ e^{-rT} (b-a)/2 Σ'_{k=0}^{N-1} A_k V_k.
\end{verbatim}

  Since
  \texttt{(b-a)/2\ ·\ A\_k\ =\ Re\{φ(kπ/(b-a);x0)\ e\^{}\{-ikπa/(b-a)\}\}},
  this is

\begin{verbatim}
P ≈ e^{-rT} Σ'_{k=0}^{N-1} Re{ φ(kπ/(b-a); x0) exp(-ikπ a/(b-a)) } V_k.
\end{verbatim}

  \textbf{Where the cf enters:} in exactly one place --- as the cosine
  coefficients of the density,
  \texttt{Re\{φ\ ·\ e\^{}\{-ikπa/(b-a)\}\}}. The spot dependence sits
  inside \(φ(·;x0)\); for BS/Lévy \(φ(u;x0)=φ_lvl(u) e^{iux0}\), so one
  set of \({V_k}\) prices the whole spot/strike grid (only the cheap
  \(e^{iux0}\) changes). Cost \(O(N)\), no integration, no FFT.
\item
  How do you pick {[}a, b{]}? What breaks at short maturity, and how did
  you fix it in your assignment?

  \textbf{A.} \textbf{Cumulant rule (Fang--Oosterlee):}

\begin{verbatim}
[a,b] = [ c1 − L√(c2+√c4),  c1 + L√(c2+√c4) ],   L ∈ [8,12] (typ. 10),
\end{verbatim}

  where \(c1\) = mean (centres the box), \(c2\) = variance (\(√c2\) is
  the natural width scale, so \(L√c2\) is an ``L-sigma'' box), \(√c4\)
  (fourth cumulant) fattens the box for heavy tails / kurtosis (zero for
  Gaussian), and \(L\) is the safety multiplier --- larger \(L\) lowers
  the truncation floor but needs more terms.

  \textbf{BS cumulants} of \texttt{ln(S\_T/K)}:

\begin{verbatim}
c1 = ln(S0/K) + (r − ½σ²)T,   c2 = σ²T,   c3 = c4 = 0   ⇒   [a,b] = c1 ± Lσ√T.
\end{verbatim}

  \textbf{What breaks as T→0:} \(c2 = σ²T → 0\), so the half-width
  \(L√(c2+√c4) → 0\) and the box collapses toward the single point
  \(c1\). With \(c4=0\) (Gaussian) there is no \(√c4\) floor to keep it
  open. Two failure modes: (i) the box gets too narrow to cover where
  \(V·f\) has mass; (ii) when pricing a whole grid of spots/strikes, a
  range centred by \(c1\) for one \((S0,K)\) is centred away from other
  grid points, so the cosine reconstruction is evaluated where it is
  meaningless --- large, erratic errors at short maturity.

  \textbf{Diagnostic signature} (this is the tell): the error
  \emph{grows} as \(T→0\) and does \emph{not} improve when \(N\)
  increases. That distinguishes a \emph{range} error (a fixed accuracy
  floor) from a \emph{series} error (curable by more terms).

  \textbf{Assignment-2 fix:} the truncation range was too narrow at
  short maturities because \(c2\) was tiny. I (i) raised \(L\) at short
  \(T\), (ii) imposed a minimum half-width floor so \(b−a\) cannot
  collapse (equivalently kept a \(√c4\)-type term), and (iii) made sure
  the range was computed to cover the entire spot/strike grid actually
  being priced, not a single point. After that the error stopped growing
  as \(T→0\) and reverted to exponential decay in \(N\).
\item
  Why cosine series instead of a complex Fourier series?

  \textbf{A.} Three reasons:

  \begin{itemize}
  \tightlist
  \item
    \textbf{(i) Real coefficients / real arithmetic.} The even extension
    of \(f\) on \([a,b]\) kills every sine term, so the surviving
    \(A_k\) are real (\(F_k\) is literally a \(Re{·}\)). A full complex
    series carries twice the coefficients and complex bookkeeping for no
    gain.
  \item
    \textbf{(ii) Direct sampling from φ.} The cosine coefficient
    \emph{is} the real part of \(φ\) evaluated at the grid frequency
    \texttt{kπ/(b-a)}:
    \texttt{A\_k\ =\ (2/(b-a))\ Re\{φ(kπ/(b-a))\ e\^{}\{-ikπa/(b-a)\}\}}.
    The cosine basis is exactly the one whose coefficients are read off
    the cf in one line --- no inversion integral.
  \item
    \textbf{(iii) Payoff side closes in elementary functions.} \(V_k\)
    are integrals of (polynomial/exponential)×cos, which integrate to
    the closed forms \texttt{χ\_k,\ ψ\_k}. With a complex-exponential
    basis the call coefficients are not as clean. The whole method pairs
    the cosine coefficients of \(f\) (from \(φ\)) against those of \(V\)
    (closed form).
  \end{itemize}
\item
  What sets the convergence rate? Give a case where it is only
  algebraic.

  \textbf{A.} Three error sources: (1) \textbf{range truncation}
  \([a,b]\) --- \emph{fixed} once chosen, it sets the accuracy
  \emph{floor} at which convergence saturates (exponentially small in
  the half-width for exponentially-decaying tails); (2) \textbf{series
  truncation} at \(N\) --- governed by the decay of the cosine
  coefficients \(A_k\); (3) \textbf{cf/model error} --- feeds straight
  through.

  The rate in \(N\) is set by the \textbf{smoothness of \(f\)}:

  \begin{itemize}
  \tightlist
  \item
    \(f\) analytic in a strip / \(f∈C^∞\) with integrable derivatives ⇒
    \(A_k\) decay \textbf{exponentially} ⇒ \textbf{exponential}
    convergence in \(N\) (machine precision in dozens--low hundreds of
    terms).
  \item
    \(f\) with only \(β\) integrable derivatives ⇒
    \(A_k = O(k^{-(β+1)})\) ⇒ only \textbf{algebraic} convergence.
    \emph{Lack of smoothness is what kills the exponential rate.}
  \end{itemize}

  \textbf{Algebraic case:} a low-regularity density ---
  e.g.~Variance-Gamma with small \(ν\) (the density develops a
  non-smooth peak / cusp), or a short-dated density approaching a Dirac,
  or near a jump/kink. There \(A_k\) decay only polynomially and you
  need many more terms.
\item
  Write the Heston characteristic function you used as ground truth.
  State the Feller condition.

  \textbf{A.} Heston: \(dS = rS dt + √v S dW1\),
  \(dv = κ( − v)dt + η√v dW2\), \(d⟨W1,W2⟩ = ρ dt\), with \(κ\)
  mean-reversion speed, \(\) long-run variance, \(η\) vol-of-vol, \(ρ\)
  correlation, \(v0\) initial variance. The characteristic function of
  the log-price \texttt{x\_T\ =\ ln(S\_T/S0)} (or \(ln S_T\)) is
  affine-exponential:

\begin{verbatim}
φ(u;T) = exp( i u (ln S0 + rT)
            + (v0/η²) · ( (1 − e^{-DT})/(1 − G e^{-DT}) ) · (κ − iρηu − D)
            + (κ /η²) · ( (κ − iρηu − D)T − 2 ln( (1 − G e^{-DT})/(1 − G) ) ) ),
\end{verbatim}

  with

\begin{verbatim}
D = √( (κ − iρηu)² + (u² + iu) η² ),
G = (κ − iρηu − D) / (κ − iρηu + D).
\end{verbatim}

  \(D\) and \(G\) come from solving the Riccati ODE for the variance
  factor of the affine system. (Equivalent ``little-trap'' form uses
  \(G\) rather than \texttt{g=1/G} to keep the complex log on the
  principal branch and avoid branch-cut discontinuities --- the cf must
  be evaluated continuously or COS picks up \(φ\)-error that feeds
  straight through.)

  \textbf{Feller condition:} \(2κ ≥ η²\). It guarantees the CIR
  variance process stays strictly positive (never hits 0), which keeps
  the density well-behaved and \(φ\) clean; if violated, \(v\) can touch
  zero and the simulation/quadrature near \(v=0\) needs care.
\item
  How does COS extend to Barrier options? To American options?

  \textbf{A. Barrier (Lecture 10), discretely monitored on
  \texttt{t\_1\textless{}…\textless{}t\_M=T}.} Knock-out survival is a
  product of indicators \texttt{∏\_m\ 1\_\{x\_m\textless{}h\}}
  (\texttt{h\ =\ ln(H/K)}). By the \textbf{Markov property} the joint
  density factorises into one-step transitions, so the \(M\)-dimensional
  integral nests into a \textbf{backward recursion}. Define continuation
  and option value at each date:

\begin{verbatim}
c(x,t_{m-1}) = e^{-rΔt} ∫ v(y,t_m) f(y|x) dy,
v(x,t_{m-1}) = R (knocked out) for x ≥ h,   = c(x,t_{m-1}) (alive) for x < h.
\end{verbatim}

  Each continuation integral is \emph{exactly} the European COS sum:

\begin{verbatim}
ĉ(x,t_{m-1}) = e^{-rΔt} Σ'_k F_k(x) V_k(t_m),   F_k(x)=Re{φ(kπ/(b-a);x) e^{-ikπa/(b-a)}}.
\end{verbatim}

  The new ingredient: the ``payoff'' \(v(y,t_m)\) is produced by the
  previous step, so its coefficients \(V_k\) are \textbf{propagated by
  backward induction}. Because \(v\) is piecewise at \(h\), \(V_k\)
  splits at the barrier:

\begin{verbatim}
_k(t_m) = Ĉ_k(a,h,t_m) + e^{-r(T-t_{m-1})} (2R/(b-a)) Ψ_k(h,b),
\end{verbatim}

  and \(Ĉ_k\) (inserting the COS continuation) becomes
  \texttt{Σ\_j\ φ\_lvl(jπ/(b-a))\ V\_j\ M\_\{k,j\}} with an analytic
  coupling \(M_{k,j}\). \(M_{k,j}\) splits into a \textbf{Hankel} part
  (\(j+k\)) and a \textbf{Toeplitz} part (\(j−k\));
  Toeplitz/Hankel×vector = convolution = 3 FFTs by circulant embedding ⇒
  \(O(N log N)\) per step, total \(O((M−1)N log₂N)\). So --- unlike
  European COS --- the barrier method \emph{does} use the FFT, because
  it must move the whole coefficient vector at every step. ``In''
  barriers come from in--out parity. Continuous monitoring is the
  \(M→∞\) limit, recovered by Richardson extrapolation in \(Δt\) using
  the known \(O(√Δt)\) discrete-vs-continuous bias. Heston ⇒ 2-D state
  (log-stock, variance): 1-D COS in log-stock + numerical integration
  over variance, or full 2-D COS; Feller \(2κ≥η²\).

  \textbf{B. American (Lecture 9).} Same backward-recursion idea but the
  decision is \emph{early exercise} rather than a barrier knock-out. On
  a time grid, at each step the option value is
  \texttt{v(x,t\_m)\ =\ max(\ g(x),\ c(x,t\_m)\ )} --- the larger of
  immediate exercise payoff \(g\) and the COS continuation value
  \texttt{c(x,t\_m)=e\^{}\{-rΔt\}Σ\textquotesingle{}\_k\ F\_k(x)V\_k(t\_\{m+1\})}.
  The early-exercise boundary \(x*_m\) is the point where \(g = c\);
  below/above it (call vs put) the option is exercised. So each step (i)
  finds \(x*_m\) (a 1-D root-find on \(g − c = 0\)), (ii) splits the
  cosine coefficients of \(v\) at \(x*_m\) into a continuation region
  (cosine coeffs of \(c\), computed from the previous \(V_k\) via the
  same analytic coupling integrals) and an exercise region (cosine
  coeffs of \(g\), closed-form \texttt{χ\_k/ψ\_k}). This is the COS
  Bermudan pricer; the American price is obtained by Richardson
  extrapolation in the number of exercise dates. Structurally identical
  to the barrier recursion, with the moving boundary \(x*_m\) playing
  the role the fixed barrier \(h\) plays there.
\item
  Compare COS with Carr--Madan. Advantages and limitations.

  \textbf{A.} Both need only the characteristic function. Carr--Madan
  damps the call (\texttt{c\_T(k)=e\^{}\{αk\}C\_T(k)}, \(α>0\), to make
  it \(L¹∩L²\)), Fourier-transforms it in closed form in terms of \(φ\),
  and recovers the price by an \textbf{inverse FFT} --- i.e.~a
  trapezoidal-type quadrature of an oscillatory integrand,
  \(O(N log N)\).

  \begin{longtable}[]{@{}
    >{\raggedright\arraybackslash}p{(\columnwidth - 4\tabcolsep) * \real{0.3333}}
    >{\raggedright\arraybackslash}p{(\columnwidth - 4\tabcolsep) * \real{0.3333}}
    >{\raggedright\arraybackslash}p{(\columnwidth - 4\tabcolsep) * \real{0.3333}}@{}}
  \toprule\noalign{}
  \begin{minipage}[b]{\linewidth}\raggedright
  \end{minipage} & \begin{minipage}[b]{\linewidth}\raggedright
  COS
  \end{minipage} & \begin{minipage}[b]{\linewidth}\raggedright
  Carr--Madan / FFT
  \end{minipage} \\
  \midrule\noalign{}
  \endhead
  \bottomrule\noalign{}
  \endlastfoot
  Core op & one \(O(N)\) cosine sum, no integration & FFT quadrature,
  \(O(N log N)\) \\
  Free params & none (range from cumulants) & damping \(α\) must be
  tuned \\
  Convergence (smooth \(f\)) & exponential in \(N\) & algebraic
  (quadrature-limited) \\
  Terms for \textasciitilde machine prec. & dozens--hundreds & thousands
  of nodes \\
  Greeks & analytic, same sum & extra transforms / finite differences \\
  \end{longtable}

  \textbf{COS wins} for smooth/piecewise-smooth densities with known
  \(φ\) and exponentially-decaying tails --- essentially all European
  pricing under BS, Heston, and standard Lévy/AJD models --- because it
  converges exponentially, needs no free parameter, and gives Greeks for
  free.

  \textbf{Limitations / be careful:} genuinely low-regularity densities
  (slow algebraic \(A_k\) decay), very short maturities, or fat tails,
  where \([a,b]\) must be chosen carefully (Q2) and \(N\) raised; and
  any error in \(φ\) (e.g.~a mis-evaluated Heston cf across the branch
  cut) feeds straight through. Carr--Madan's FFT also naturally returns
  a \emph{whole strip of strikes} in one transform, which COS matches
  only because its \({V_k}\) are strike-independent in log-moneyness.
\end{enumerate}

\hypertarget{my-answers-concise-board-versions-re-derive-these-out-loud-dont-read}{%
\section{My answers (concise board versions --- re-derive these out
loud, don't
read)}\label{my-answers-concise-board-versions-re-derive-these-out-loud-dont-read}}

Compressed versions of the full answers above; expand any of them on
demand.

\begin{enumerate}
\def\labelenumi{\arabic{enumi}.}
\item
  \textbf{COS call price.} \(P = e^{-rT}∫V f dy\); I don't know \(f\)
  but know its cf \(φ\). Truncate to \([a,b]\), cosine-expand \(f\), and
  the coefficients come straight off \(φ\):
  \texttt{A\_k\ ≈\ (2/(b-a))Re\{φ(kπ/(b-a))e\^{}\{-ikπa/(b-a)\}\}}
  (extend the coefficient integral to \(R\) --- the only approximation).
  Pair against payoff coeffs
  \texttt{V\_k=(2/(b-a))∫\_a\^{}b\ V\ cos(...)dy}; the integral of the
  product is \texttt{(b-a)/2\ Σ\textquotesingle{}A\_kV\_k}, giving
  \texttt{P\ ≈\ e\^{}\{-rT\}\ Σ\textquotesingle{}\_\{k=0\}\^{}\{N-1\}\ Re\{φ(kπ/(b-a);x0)e\^{}\{-ikπa/(b-a)\}\}\ V\_k}.
  The cf enters only as the cosine coefficients of the density; \(x0\)
  lives inside \(φ\), so one \({V_k}\) prices the whole grid.
\item
  \textbf{\([a,b]\).} \(c1 ± L√(c2+√c4)\), \(L≈10\): \(c1\) centres,
  \(√c2\) is the width scale, \(√c4\) fattens for fat tails, \(L\) is
  the safety margin. BS: \texttt{c1=ln(S0/K)+(r-½σ²)T}, \(c2=σ²T\),
  \(c4=0\), so \([a,b]=c1±Lσ√T\). As \(T→0\), \(c2→0\) so the box
  collapses (no \(√c4\) floor in the Gaussian case) → error \emph{grows}
  as \(T→0\) and is \emph{not} cured by more \(N\) (that signature =
  range, not series, error). Assignment-2 fix: raise \(L\) at short
  \(T\), impose a minimum half-width floor, and size the range to the
  whole spot/strike grid, not one point.
\item
  \textbf{Why cosine.} Even extension ⇒ real \(A_k\), real arithmetic,
  half the storage; \(A_k\) is literally \(Re{φ}\) at the grid frequency
  (direct cf sampling); payoff coeffs reduce to elementary
  \texttt{χ\_k,\ ψ\_k}. A complex series gives none of these cleanly.
\item
  \textbf{Convergence.} Three errors: range (fixed floor), series (decay
  of \(A_k\)), cf/model (passes through). Rate set by smoothness of
  \(f\): analytic ⇒ exponential \(A_k\) decay ⇒ exponential in \(N\);
  only \(β\) integrable derivatives ⇒ \(A_k=O(k^{-(β+1)})\) ⇒ algebraic.
  Algebraic case: VG with small \(ν\), or short-dated density near a
  Dirac/kink.
\item
  \textbf{Heston cf + Feller.} Affine-exponential cf in \(u\) with
  \(D=√((κ-iρηu)²+(u²+iu)η²)\), \texttt{G=(κ-iρηu-D)/(κ-iρηu+D)},
  variance term
  \texttt{(v0/η²)((1-e\^{}\{-DT\})/(1-Ge\^{}\{-DT\}))(κ-iρηu-D)} plus
  the \texttt{(κ/η²)} drift term; use the ``little-trap'' form for a
  continuous branch. Feller: \(2κ ≥ η²\) keeps \(v>0\) and \(φ\) clean.
\item
  \textbf{Barrier / American.} Both = backward recursion via Markov
  factorisation; each step is European COS
  \texttt{ĉ=e\^{}\{-rΔt\}Σ\textquotesingle{}F\_k(x)V\_k(t\_m)} with
  \(V_k\) propagated. Barrier: knock out by overwriting \(v=R\) above
  \(h\), \(V_k\) splits at \(h\), the coupling \(M_{k,j}\) is
  Hankel+Toeplitz ⇒ FFT, \(O((M-1)N log₂N)\). American: \(v=max(g,c)\),
  split \(V_k\) at the early-exercise boundary \(x*_m\) (root of
  \(g-c\)); Bermudan + Richardson in \#exercise dates. The moving
  boundary \(x*\) plays the role of the fixed barrier \(h\).
\item
  \textbf{COS vs Carr--Madan.} Both need only \(φ\). CM damps the call
  (param \(α\)) and inverts by FFT quadrature --- algebraic,
  \(O(N log N)\), tunable \(α\). COS is a single \(O(N)\) cosine sum, no
  integration, no free parameter, exponential convergence, free Greeks.
  COS wins for smooth densities with decaying tails; be careful for
  low-regularity / very short \(T\) / fat tails (set \([a,b]\), raise
  \(N\)), and any \(φ\) error feeds straight through.
\end{enumerate}

\begin{center}\rule{0.5\linewidth}{0.5pt}\end{center}

Lecture 6 drill set (European COS)

Companion to \texttt{notes/Lecture06-notes.pdf}. Each question has a
terse answer key in the collapsible note --- cover it and answer out
loud first. ``Derive'' means on paper, no notes.

\hypertarget{a.-density-recovery}{%
\section{A. Density recovery}\label{a.-density-recovery}}

\textbf{A1.} Write the half-range cosine expansion of a density f on
{[}a,b{]}, with the prime convention. What does the prime mean and why
is A0/2 the mean of f over {[}a,b{]}? \textgreater{} Key: f(x)=Σ' A\_k
cos(kπ(x-a)/(b-a)), A\_k=(2/(b-a))∫\_a\^{}b f cos(kπ(x-a)/(b-a))dx.
Prime halves k=0. A0 = (2/(b-a))∫f = (2/(b-a))·(mass), and A0/2 over the
constant basis function is the average height.

\textbf{A2.} \emph{(Core derivation.)} Derive A\_k ≈ (2/(b-a)) Re\{
φ(kπ/(b-a)) exp(-ikπa/(b-a)) \} starting from the coefficient integral.
Identify the single approximation you make and where. \textgreater{}
Key: cos = Re of exp; pull out e\^{}\{-ikπa/(b-a)\}; the remaining
∫\_a\^{}b f e\^{}\{i(kπ/(b-a))x\}dx is extended to ℝ (THE approximation
--- range truncation) = φ(kπ/(b-a)). Everything else exact.

\textbf{A3.} Standard normal: φ(t)=e\^{}\{-t²/2\}. Write F\_k for a=-b
symmetric. What range do the slides suggest and why (10\^{}\{-8\} clip)?
\textgreater{} Key: substitute φ(t)=e\^{}\{-t²/2\} at t=kπ/(b-a) into
F\_k, giving F\_k=(2/(b-a))Re\{e\^{}\{-½(kπ/(b-a))²\}
e\^{}\{-ikπa/(b-a)\}\}. The slides set a=Φ\^{}\{-1\}(1e-8), b=-a,
i.e.~clip each tail at probability 1e-8 --- symmetric because the
standard normal is symmetric about 0. The Gaussian factor
e\^{}\{-½(kπ/(b-a))²\} makes F\_k decay super-exponentially in k, which
is the canonical exponential-convergence picture.

\hypertarget{b.-why-cosine}{%
\section{B. Why cosine}\label{b.-why-cosine}}

\textbf{B1.} Three reasons COS uses a cosine series, not a complex
Fourier series. Be specific about each. \textgreater{} Key: (i) even
extension ⇒ real coefficients, real arithmetic, half the storage; (ii)
A\_k is literally Re\{φ at grid frequency\} --- direct sampling from
ch.f.; (iii) payoff coefficients V\_k reduce to elementary χ\_k, ψ\_k
(exp/poly × cos integrate in closed form).

\hypertarget{c.-pricing-formula}{%
\section{C. Pricing formula}\label{c.-pricing-formula}}

\textbf{C1.} \emph{(Core derivation.)} From P = e\^{}\{-rT\}∫ V f dy,
derive P ≈ e\^{}\{-rT\} Σ'\_\{k=0\}\^{}\{N-1\} Re\{φ(kπ/(b-a);x0)
e\^{}\{-ikπa/(b-a)\}\} V\_k. State the V\_k definition and the
sum--integral swap justification. \textgreater{} Key: insert cosine
expansion of f, swap Σ/∫ (uniform convergence, smooth f), define
V\_k=(2/(b-a))∫\_a\^{}b V cos(\ldots)dy ⇒ P≈e\^{}\{-rT\}(b-a)/2 Σ' A\_k
V\_k; sub (b-a)/2·A\_k = Re\{φ e\^{}\{-ikπa/(b-a)\}\}; truncate at N.

\textbf{C2.} Where exactly does the ch.f. enter, and where does the spot
x0 dependence live? Why does one set of \{V\_k\} price the whole
strike/spot grid? \textgreater{} Key: φ enters only as the cosine
coefficients of the density, Re\{φ e\^{}\{-ikπa/(b-a)\}\} --- nowhere
else. The spot x0 sits inside φ(·;x0); for Lévy/BS
φ(u;x0)=φ\_lvl(u)e\^{}\{iux0\}, so across a strike/spot grid only the
cheap e\^{}\{iux0\} phase changes while the level cf φ\_lvl is reused.
The payoff coeffs V\_k are strike-independent in the log-moneyness
variable, so they are computed once and reused for the whole grid ---
that is the source of COS's grid efficiency.

\textbf{C3.} What is the per-price computational cost of the COS sum?
Does it call an FFT? \textgreater{} Key: O(N) per price --- a single
finite cosine sum of N terms, no numerical integration and no FFT. Each
term is one cf evaluation times a precomputed V\_k, so the whole price
is essentially a dot product. This is the headline contrast with
Carr--Madan (O(N log N) FFT quadrature) and with the discrete-barrier
COS (which does need the FFT because it propagates the full coefficient
vector at every step).

\hypertarget{d.-payoff-coefficients-v_k}{%
\section{D. Payoff coefficients V\_k}\label{d.-payoff-coefficients-v_k}}

\textbf{D1.} \emph{(Derive.)} Define χ\_k(c,d)=∫\_c\^{}d e\^{}y
cos(ω\_k(y-a))dy, ω\_k=kπ/(b-a). Derive the closed form. \textgreater{}
Key: χ\_k=Re\{e\^{}\{-iω\_k a\}/(1+iω\_k)
{[}e\^{}\{(1+iω\_k)y\}{]}\_c\^{}d\}; rationalise by (1-iω\_k)/(1+ω\_k²);
χ\_k = 1/(1+ω\_k²){[}cos(ω\_k(d-a))e\^{}d - cos(ω\_k(c-a))e\^{}c + ω\_k
sin(ω\_k(d-a))e\^{}d - ω\_k sin(ω\_k(c-a))e\^{}c{]}.

\textbf{D2.} Write ψ\_k(c,d) including the k=0 case. Why must k=0 be
special-cased? \textgreater{} Key: for k≠0,
ψ\_k(c,d)=(1/ω\_k){[}sin(ω\_k(d-a))-sin(ω\_k(c-a)){]}=(b-a)/(kπ){[}sin(ω\_k(d-a))-sin(ω\_k(c-a)){]}.
For k=0, ω\_0=0 so the integrand cos(0)=1 and ψ\_0=d-c.~The k=0 case
must be special-cased because the general 1/ω\_k formula is 0/0 there (a
removable singularity); in code a naive division by ω\_0=0 gives NaN, so
the k=0 term --- which also carries the ½ from the prime --- is set
explicitly.

\textbf{D3.} In y=ln(S\_T/K), give V\_k for call and put. Why is the
call interval {[}0,b{]} and the put interval {[}a,0{]}? \textgreater{}
Key: call V\_k=(2/(b-a))K{[}χ\_k(0,b)-ψ\_k(0,b){]} --- payoff
K(e\textsuperscript{y-1)}+ lives on y≥0; put
V\_k=(2/(b-a))K{[}ψ\_k(a,0)-χ\_k(a,0){]} --- payoff
K(1-e\textsuperscript{y)}+ lives on y≤0.

\hypertarget{e.-truncation-range-the-bug}{%
\section{E. Truncation range (THE
bug)}\label{e.-truncation-range-the-bug}}

\textbf{E1.} State {[}a,b{]}={[}c1 - L√(c2+√c4), c1 + L√(c2+√c4){]}.
Interpret each of c1, c2, √c4, L. \textgreater{} Key: c1 mean (centres),
c2 variance (√c2 width scale, L-sigma box), √c4 fattens box for heavy
tails/kurtosis (0 for Gaussian), L safety multiplier (8--12, typ. 10).

\textbf{E2.} Give the BS cumulants of ln(S\_T/K) and the resulting box.
Which cumulants vanish? \textgreater{} Key: c1=ln(S0/K)+(r-½σ²)T,
c2=σ²T, c4=0 (and c3=0) ⇒ {[}a,b{]}=c1 ± Lσ√T.

\textbf{E3.} \emph{(High priority.)} As T→0, what happens to {[}a,b{]}
and why? Give the diagnostic signature distinguishing a range problem
from a series problem. How do you fix it? \textgreater{} Key: c2=σ²T→0 ⇒
half-width→0 ⇒ box collapses to \textasciitilde c1; with c4=0 no floor
keeps it open. Signature: error GROWS as T→0 and does NOT improve when N
increases (range error is a floor, not curable by more terms). Fix:
raise L at short T; keep √c4 / impose a minimum half-width floor; ensure
the range covers the whole spot/strike grid actually priced, not one
point.

\hypertarget{f.-convergence}{%
\section{F. Convergence}\label{f.-convergence}}

\textbf{F1.} List the three error sources in the COS price and say which
one sets the accuracy floor. \textgreater{} Key: (1) range truncation
{[}a,b{]} --- FIXED once chosen, sets the floor; (2) series truncation
at N --- set by decay of A\_k; (3) ch.f./model error feeds straight
through. Convergence saturates at the range floor.

\textbf{F2.} What property of f gives exponential vs only algebraic
convergence in N? Give a case where it is only algebraic. \textgreater{}
Key: smoothness/analyticity ⇒ A\_k decay exponentially ⇒ exponential in
N. Finite smoothness (β integrable derivatives) ⇒
A\_k=O(k\^{}\{-(β+1)\}) ⇒ algebraic. Algebraic case: density with low
regularity / near a kink/jump, e.g.~VG with small ν, or short-dated
densities approaching a Dirac.

\textbf{F3.} COS vs MC convergence rate? Roughly how many terms for
machine precision on Heston (cf slide 31)? \textgreater{} Key: MC
converges at O(N\^{}\{-1/2\}) in the number of paths, whereas COS
converges exponentially in the number of terms for smooth densities. On
Heston (slide 31) N\textasciitilde160 terms reach \textasciitilde1e-10
error, versus thousands of FFT nodes for worse error and longer CPU
time. So COS reaches machine precision with N in the dozens to low
hundreds --- the gap is the whole selling point of the method.

\hypertarget{g.-greeks}{%
\section{G. Greeks}\label{g.-greeks}}

\textbf{G1.} Why are Delta and Vega ``free'' in COS? Write the Delta
sum. \textgreater{} Key: S0 (via x=ln(S0/K)) and the model parameters
enter the COS sum only analytically, through φ(·;x0)=φ(·)e\^{}\{iux\};
so differentiation passes straight through the finite sum with no
re-integration. Δ≈e\^{}\{-rT\}Σ' Re\{φ(kπ/(b-a)) e\^{}\{ikπ(x-a)/(b-a)\}
· ikπ/(b-a)\} V\_k/S0 (chain rule ∂x/∂S0=1/S0). Vega differentiates φ
w.r.t. the variance parameter instead. The same \{V\_k\} are reused, so
Greeks cost essentially nothing extra --- a structural advantage over
quadrature/FFT pricers that need finite differences.

\begin{center}\rule{0.5\linewidth}{0.5pt}\end{center}

Lecture 10 drill set (COS for barrier options)

Companion to \texttt{notes/Lecture10-notes.pdf}. Same cover-and-answer
format.

\hypertarget{h.-the-pricing-problem}{%
\section{H. The pricing problem}\label{h.-the-pricing-problem}}

\textbf{H1.} Define up-and-out vs up-and-in call. Why are there exactly
8 barrier types? \textgreater{} Key: up-and-out call pays max(S\_T-K,0)
UNLESS the asset crosses B\textgreater S0 during {[}0,T{]}, in which
case it pays 0 (knocked out); up-and-in call pays 0 UNLESS B is crossed,
in which case it becomes a vanilla call (knocked in). The 8 types =
up/down × in/out × call/put (2×2×2). Each has a closed-form GBM price.

\textbf{H2.} State in--out parity and give its two practical uses.
\textgreater{} Key: holding knock-in and knock-out with the same (K,B,T)
reproduces the vanilla, because exactly one of the two is alive at
maturity: C\_in + C\_out = e\^{}\{-r(T-t)\}E{[}(S\_T-K)\^{}+{]}. Use
(i): price the easier leg (usually the knock-out) plus the vanilla, then
back out the other by subtraction --- cheaper and more stable than
handling the in-condition directly. Use (ii): it gives a free
correctness check on any numerical barrier price.

\textbf{H3.} Write the two equivalent formulations (expectation and
PDE). Where does the barrier enter the PDE? \textgreater{} Key:
V=e\^{}\{-r(T-t)\}E\_x{[}(α(e\textsuperscript{\{X\_T\}-K))}+
1\_\{τ\_B\textgreater T\}{]}; or localized Feynman--Kac BS PDE on
{[}a,b{]} with terminal payoff and ABSORBING boundary V(a,t)=V(b,t)=0.
Barrier enters via the homogeneous Dirichlet BC, not the equation.

\hypertarget{i.-reflection-principle-derive}{%
\section{I. Reflection principle
(derive)}\label{i.-reflection-principle-derive}}

\textbf{I1.} \emph{(Core.)} Derive P(τ\_a ≤ T) = 2P(W\_T ≥ a) and hence
the density of the running maximum Ξ. \textgreater{} Key: split
P(τ\_a≤T) by \{W\_T\textgreater a\},\{W\_T≤a\}; reflection ⇒ both halves
equal ⇒ =2P(τ\_a≤T,W\_T≥a)=2P(W\_T≥a) (since \{W\_T≥a\}⊆\{τ\_a≤T\}). So
Ξ \textasciitilde{} \textbar W\_T\textbar, P(Ξ≤a)=erf-type,
p\_Ξ(a)=√(2/πT) e\^{}\{-a²/2T\} 1\_\{a≥0\}.

\textbf{I2.} \emph{(Core.)} Derive the joint density of (Ξ\_0\^{}T,
W\_T). \textgreater{} Key: for b≤a, reflect after τ\_a:
P(Ξ≥a,W\_T\textless b)=P(Ξ≥a,W\_T\textgreater2a-b)=P(W\_T\textgreater2a-b)
(since 2a-b≥a). Differentiate -∂²/∂a∂b ⇒
p(a,b)=√(2/π)(2a-b)/T\^{}\{3/2\} e\^{}\{-(2a-b)²/2T\} 1\_\{a≥max(b,0)\}.

\textbf{I3.} How do you get from standard BM to drifted BM (GBM
log-price) results? \textgreater{} Key: Girsanov change of measure. A
drifted BM W\_T+θt under P is a driftless BM under an equivalent measure
Q, with Radon--Nikodym derivative dP/dQ = e\^{}\{θW\_T-½θ²T\}. So I
compute the reflection-principle result for driftless BM under Q, then
reweight by that exponential to recover the drifted (GBM log-price) law.
This is exactly what turns the driftless joint density of (Ξ,W\_T) into
the drifted one used in the closed-form barrier prices.

\hypertarget{j.-pde-routes}{%
\section{J. PDE routes}\label{j.-pde-routes}}

\textbf{J1.} State the method-of-images identity and the down-and-out
call it produces. What is α? \textgreater{} Key: for the BS operator
LV=V\_t+rSV\_S+½σ²S²V\_SS-rV, if LV=0 then the image
U(S,t)=S\^{}\{2α\}V(B²/S,t) also solves LU=0, with α=½(1-2r/σ²). This is
the BS analogue of reflecting a heat-equation solution across a wall.
Subtracting the image with the right weight cancels the value on the
barrier S=B, giving the down-and-out call C\_down-out=C\_van(S,K) -
(S/B)\^{}\{2α\} C\_van(B²/S,K): the first term is the vanilla, the
second removes exactly the paths that would have crossed B, evaluated at
the mirror spot B²/S.

\textbf{J2.} \emph{(Derive.)} Sine-series PDE route: what change of
variables reduces the localized BS PDE to the heat equation, and why
does the sine basis appear? \textgreater{} Key: τ=½σ²(T-t), κ=2r/σ²,
V=e\^{}\{αx+βτ\}U with α=-½(κ-1), β=-α²-κ, z=x-a ⇒ U\_τ=U\_zz on
{[}0,b-a{]}, U(0,τ)=U(b-a,τ)=0. Homogeneous Dirichlet BCs ⇒ sine series.
Propagate modes by e\^{}\{-(nπ/(b-a))²τ\}.

\hypertarget{k.-cos-for-discrete-barriers-the-main-event}{%
\section{K. COS for discrete barriers (the main
event)}\label{k.-cos-for-discrete-barriers-the-main-event}}

\textbf{K1.} \emph{(High priority.)} Why does a discretely-monitored
barrier become a BACKWARD RECURSION? Where does the Markov property
enter? \textgreater{} Key: survival = ∏\emph{\{m\}1}\{x\_m\textless h\};
Markov ⇒ joint density factorises into one-step transitions ⇒ the M-dim
integral nests ⇒ backward recursion:
c(x,t\_\{m-1\})=e\^{}\{-rΔt\}∫v(y,t\_m)f(y\textbar x)dy, then knock out:
v=R for x≥h, =c for x\textless h.

\textbf{K2.} Write the COS step for the continuation value. How is it
the European COS formula? \textgreater{} Key:
ĉ(x,t\_\{m-1\})=e\^{}\{-rΔt\}Σ'\_\{k\} F\_k(x)V\_k(t\_m),
F\_k(x)=Re{[}φ(kπ/(b-a);x)e\^{}\{-ikπa/(b-a)\}{]}, V\_k=cosine coeffs of
v(y,t\_m). It IS Euro COS with v(·,t\_m) playing the role of the payoff
--- except v's coefficients come from the previous step.

\textbf{K3.} \emph{(Core.)} Why must V\_k be propagated by backward
induction, and how does the knock-out split V\_k at h? \textgreater{}
Key: v is piecewise (R above h, ĉ below) ⇒ \emph{k(t\_m)=Ĉ\_k(a,h,t\_m)
+ e\^{}\{-r(T-t}\{m-1\})\}(2R/(b-a))Ψ\_k(h,b). Ĉ\_k inserts the COS
continuation value ⇒ Ĉ\_k=e\^{}\{-rΔt\}Re{[}Σ'\emph{j
φ\_Levy(jπ/(b-a))V\_j(t}\{m+1\})M\_\{k,j\}{]}, with M\_\{k,j\} an
analytic coupling integral.

\textbf{K4.} \emph{(High priority --- contrast with Lecture 6.)} Why
does the discrete-barrier COS use the FFT when European COS does not?
Give the matrix structure and the complexity. \textgreater{} Key:
M\_\{k,j\} splits into a Hankel part (depends on j+k) and Toeplitz part
(depends on j-k). Toeplitz/Hankel × vector = convolution = 3 FFTs via
circulant embedding ⇒ O(N log N) per step. Euro COS is a single O(N) sum
needing no coefficient propagation; the barrier needs all \{V\_k\} at
every step. Total: O((M-1)N log₂N).

\textbf{K5.} State the truncation range for the barrier problem. What
extra constraint beyond Lecture 6? \textgreater{} Key: {[}a,b{]}=(c1+x0)
± L√(c2+√c4), x0=ln(S0/K), cumulants of the one-step increment. Extra:
the barrier h=ln(H/K) must sit well INSIDE {[}a,b{]} or the split at h
is meaningless; and range is fixed once but reused for all M steps.

\textbf{K6.} How is Heston handled? Continuous monitoring?
\textgreater{} Key: Heston has a 2-D state (log-stock, variance), so the
one-step continuation integral is two-dimensional: either 1-D COS in the
log-stock with numerical integration over the variance, or a full 2-D
COS recovering the joint density. The characteristic function is
Heston's (Lecture 7) and the Feller condition 2κ≥η² keeps it
well-behaved. Continuous monitoring is the M→∞ (Δt→0) limit; in practice
the discrete price converges to it and a Richardson-type extrapolation
in Δt --- using the known O(√Δt) discrete-vs-continuous barrier bias ---
recovers the continuous price.

\end{document}
